
\documentclass{beamer}

% Configuración del tema
\usetheme{Madrid}
\usecolortheme{seahorse}
\setbeamertemplate{navigation symbols}{}
\setbeamertemplate{footline}[frame number]
\setbeamertemplate{itemize items}[circle]

% Información de la presentación
\title{Introducción y Marco Legal en PRL}
\author{Instructor: Santo Domingo Savio}
\date{\today}

\begin{document}

% Portada
\begin{frame}
    \titlepage
    \note{
        Bienvenida a los participantes. Introducir el objetivo de la sesión: aprender los fundamentos de la PRL y su relevancia en el ámbito laboral.
    }
\end{frame}

% Índice
\begin{frame}{Agenda}
    \tableofcontents
    \note{
        Explicar brevemente los temas que se cubrirán, destacando la importancia de cada uno para entender la PRL.
    }
\end{frame}

\section{M1: Introducción a la PRL}
% Sección 1: Introducción a la PRL
\begin{frame}{¿Qué es la PRL?}
    \begin{itemize}
        \item Conjunto de medidas destinadas a evitar riesgos para la salud y seguridad en el entorno laboral.
        \item Prevención vs. reacción: una inversión en seguridad.
        \item Rol clave en todos los sectores.
    \end{itemize}
    \note{
        Ejemplo: Comparar la PRL con un seguro de coche: es mejor prevenir un accidente que lidiar con sus consecuencias.
    }
\end{frame}

% Sección 2: Conceptos básicos de PRL
\begin{frame}{Definiciones clave}
    \begin{itemize}
        \item \textbf{Riesgo:} Probabilidad de que ocurra un daño.
        \item \textbf{Peligro:} Fuente potencial de daño.
        \item \textbf{Accidente:} Evento que genera un daño.
    \end{itemize}
    \note{
        Ejemplo: Diferenciar entre riesgo y peligro con una cafetera caliente: el peligro es la cafetera misma; el riesgo es quemarse al tocarla.
    }
\end{frame}

\begin{frame}{Tipos de riesgos}
    \begin{itemize}
        \item Físicos
        \item Químicos
        \item Ergonómicos
        \item Psicosociales
        \item Eléctricos
    \end{itemize}
    \note{
        Ejemplo: En una oficina, un riesgo ergonómico común es una silla mal ajustada, mientras que un riesgo eléctrico puede ser un cable expuesto.
    }
\end{frame}

% Sección 3: Importancia en el ámbito 
\begin{frame}{Repercusión en productividad y salud}
    \begin{itemize}
        \item Manipulación de hardware y componentes electrónicos.
        \item Exposición a radiación electromagnética.
        \item Accidentes comunes: descargas eléctricas, posturas inadecuadas.
    \end{itemize}
    \note{
        Anécdota: Un trabajador se lesionó al levantar incorrectamente un servidor pesado. Esto subraya la importancia de la capacitación en técnicas de levantamiento seguro.
    }
\end{frame}

% Sección 4: Normativa y legislación básica
\begin{frame}{Ley y reglamentos principales}
    \begin{itemize}
        \item Ley 31/1995 de Prevención de Riesgos Laborales.
        \item Reglamento RD 39/1997 y otras normativas clave.
    \end{itemize}
    \note{
        Ejemplo: Mencionar un caso en el que una empresa fue sancionada por no cumplir con la evaluación de riesgos.
    }
\end{frame}

\begin{frame}{Obligaciones de empresa y trabajador}
    \textbf{Obligaciones de la empresa:}
    \begin{itemize}
        \item Evaluación de riesgos.
        \item Formación continua.
        \item Vigilancia de la salud.
    \end{itemize}
    \textbf{Obligaciones del trabajador:}
    \begin{itemize}
        \item Uso correcto de EPI.
        \item Cumplimiento de protocolos.
    \end{itemize}
    \note{
        Ejemplo: La empresa debe proporcionar cascos de seguridad, pero el trabajador también debe usarlos correctamente.
    }
\end{frame}

% Sección 5: Organización de la prevención
\begin{frame}{Servicios de prevención}
    \begin{itemize}
        \item Propios
        \item Ajenos
        \item Mancomunados
    \end{itemize}
    \note{
        Explicar las ventajas y desventajas de cada tipo de servicio con ejemplos reales.
    }
\end{frame}

\begin{frame}{Documentos clave}
    \begin{itemize}
        \item Evaluación de riesgos.
        \item Plan de emergencia.
    \end{itemize}
    \note{
        Ejemplo: Mostrar un fragmento de un plan de emergencia típico para un edificio de oficinas.
    }
\end{frame}

% Sección 6: Derechos y obligaciones
\begin{frame}{Derechos del trabajador}
    \begin{itemize}
        \item Información y formación.
        \item Participación en comités de seguridad.
    \end{itemize}
    \note{
        Ejemplo: Los trabajadores tienen derecho a recibir capacitación específica si hay nuevos riesgos en el lugar de trabajo, como maquinaria nueva.
    }
\end{frame}

\begin{frame}{Obligaciones del trabajador}
    \begin{itemize}
        \item Uso adecuado de los equipos de protección.
        \item Notificación de riesgos y accidentes.
    \end{itemize}
    \note{
        Anécdota: Un empleado detectó un cable suelto y lo notificó a tiempo, evitando un accidente.
    }
\end{frame}


\section{M2: Identificación de Riesgos y Peligrosidad}

% Sección: Introducción
\note{
Ejemplo para destacar: En 2019, una empresa de tecnología sufrió pérdidas millonarias porque un técnico no utilizó una pulsera antiestática y dañó componentes clave durante una reparación.
}

% Sección: Clasificación de los riesgos

\begin{frame}{Clasificación de los riesgos}
    \begin{itemize}
        \item \textbf{Riesgos físicos:} electricidad estática, temperaturas elevadas, radiaciones, ruido.
        \item \textbf{Riesgos eléctricos:} descargas eléctricas, cables defectuosos, sobrecargas.
        \item \textbf{Riesgos mecánicos:} caídas, golpes, cortes.
        \item \textbf{Riesgos ergonómicos:} posturas forzadas, manipulación de cargas, tareas repetitivas.
        \item \textbf{Riesgos químicos:} vapores tóxicos, baterías peligrosas.
        \item \textbf{Riesgos psicosociales:} estrés laboral, trabajo a turnos.
    \end{itemize}
\end{frame}

\note{
Anecdota: En una sala de servidores mal ventilada, las temperaturas elevadas causaron una falla en varios equipos, interrumpiendo operaciones críticas por horas.
}

% Sección: Evaluación de la peligrosidad

\begin{frame}{Evaluación de la Peligrosidad}
    \begin{itemize}
        \item \textbf{Probabilidad de ocurrencia:} frecuencia con la que puede suceder el riesgo.
        \item \textbf{Severidad de las consecuencias:} impacto en caso de que ocurra.
    \end{itemize}
    \vspace{0.5cm}
    \textbf{Ejemplo:} Trabajar sin pulsera antiestática puede dañar componentes, lo que genera costos elevados de reparación.
    \paragraph{Métodos de Identificación de Riesgos}
    \begin{itemize}
        \item Inspecciones visuales.
        \item Revisión de protocolos.
        \item Uso de checklists de seguridad.
    \end{itemize}
\end{frame}

\note{
Usa una matriz de riesgos para ayudar a visualizar las áreas críticas. Explica cómo una sobrecarga eléctrica podría no ser frecuente, pero sus consecuencias son severas.
}

% Sección: Identificación de riesgos

\begin{frame}{Métodos de Identificación de Riesgos}
    \begin{itemize}
        \item Inspecciones visuales.
        \item Revisión de protocolos.
        \item Uso de checklists de seguridad.
    \end{itemize}
\end{frame}

\note{
Ejemplo real: Un técnico evitó un incendio gracias a la inspección visual que identificó un cable con aislamiento dañado.
}

% Sección: Casos prácticos

\begin{frame}{Actividad: Identificación de Riesgos}
    \textbf{Instrucciones:}
    \begin{enumerate}
        \item Observa las imágenes o diagramas presentados.
        \item Identifica situaciones inseguras.
        \item Discute en grupo los posibles accidentes.
    \end{enumerate}
\end{frame}

\note{
Imágenes sugeridas: Una sala de servidores con cables desordenados y un técnico sin equipo de protección. Promueve la discusión sobre los riesgos visibles y ocultos.
}

%\note{ Cierra destacando cómo prevenir riesgos no solo protege equipos y recursos, sino también a las personas, que son el recurso más valioso.}

\section{M3: Uso Seguro de Máquinas y Herramientas}

% Secciones principales
\begin{frame}{Importancia de la Seguridad}
    \begin{itemize}
        \item La seguridad garantiza la integridad física y el buen uso del equipamiento.
        \item Previene accidentes y reduce costos por reparaciones o reposiciones.
        \item Ejemplo: En 2022, una empresa tecnológica reportó pérdidas millonarias debido a incendios causados por soldadores mal manejados.
    \end{itemize}
    \note{
        Comienza destacando la importancia de la prevención. Usa el ejemplo del incendio para captar la atención y generar reflexión.}{}
\end{frame}

\begin{frame}{Herramientas Manuales}
    \begin{itemize}
        \item \textbf{Ejemplos}: destornilladores, llaves, crimpadoras, alicates, pinzas.
        \item Uso principal: ensamblaje y ajustes precisos.
        \item Precaución: comprueba el aislamiento y el estado de las puntas.
    \end{itemize}
    \includegraphics[width=0.5\textwidth]{herramientas.jpg}
    %Comenta cómo el desgaste de una herramienta manual, como un destornillador, puede provocar deslizamientos y accidentes. Anécdota: Un técnico sufrió un corte profundo porque no revisó su herramienta.
\end{frame}

\begin{frame}{Herramientas Eléctricas y Auxiliares}
    \begin{itemize}
        \item \textbf{Ejemplos}: taladros, soldadores, pistolas de aire caliente, multímetros.
        \item Equipos auxiliares: bancadas, soportes, ventiladores.
        \item Recomendación: Desconectar antes de cambiar piezas.
    \end{itemize}
    \includegraphics[width=0.5\textwidth]{equipos.jpg}
    %Destaca la importancia de los equipos auxiliares para la estabilidad. Ejemplo: Una bancada mal ajustada puede generar errores en soldaduras, arruinando el trabajo.
\end{frame}


\begin{frame}{Normas Generales}
    \begin{itemize}
        \item Revisión previa: cables, enchufes y desgaste de piezas.
        \item Uso adecuado: técnicas correctas para evitar lesiones.
        \item Desconexión y bloqueo antes de mantenimiento.
    \end{itemize}
    \includegraphics[width=0.5\textwidth]{seguridad.jpg}
   %Comparte un caso donde un técnico olvidó desconectar un soldador, causando un incendio menor. Resalta que las revisiones previas ahorran problemas.
\end{frame}

\begin{frame}{Principales Causas}
    \begin{itemize}
        \item Uso indebido o apresurado.
        \item Falta de formación.
        \item No usar EPI o ignorar la señalización.
        \item Mantenimiento deficiente.
    \end{itemize}
    \includegraphics[width=0.5\textwidth]{accidentes.jpg}
    %Relata cómo un multímetro mal calibrado causó un corto circuito en un taller, provocando pérdida de datos importantes.
\end{frame}


\begin{frame}{Pasos Clave}
    \begin{itemize}
        \item Apagar máquinas eléctricas de forma segura.
        \item Uso correcto de extintores y sistemas de detección.
        \item Seguir el plan de evacuación.
    \end{itemize}
    \includegraphics[width=0.5\textwidth]{emergencias.jpg}
 %Explica el uso de extintores y cómo las alarmas son cruciales. Ejemplo: un caso real donde una alarma evitó un desastre mayor en un centro de datos.}{}
\end{frame}


\begin{frame}{Características de Seguridad}
    \begin{itemize}
        \item Protecciones y carcasas.
        \item Alarmas e indicadores luminosos.
        \item Recomendaciones del fabricante.
    \end{itemize}
    \includegraphics[width=0.5\textwidth]{seguridad_maquinas.jpg}
%Habla de las recomendaciones del fabricante y cómo ignorarlas puede invalidar garantías o provocar accidentes evitables.}{}
\end{frame}


\section{M4: Equipos de Protecci\'on y Medidas de Seguridad} 
% Contenidos: Tipos de EPI
\begin{frame}{Tipos de EPI y su utilidad}
    \begin{itemize}
        \item \textbf{Protecci\'on de la cabeza:} Cascos para \'areas con riesgo de golpes o ca\'idas de objetos.
        \item \textbf{Protecci\'on ocular:} Gafas de seguridad o pantallas faciales para soldadura o trabajos con chispas.
        \item \textbf{Protecci\'on de manos:} Guantes anticortes, guantes antiest\'aticos.
        \item \textbf{Protecci\'on de pies:} Calzado de seguridad con puntera reforzada o protecci\'on electrost\'atica (ESD).
        \item \textbf{Protecci\'on respiratoria:} Mascarillas y filtros en tareas de limpieza o uso de productos qu\'imicos.
    \end{itemize}
\end{frame}

% Contenidos: Normativa
\begin{frame}{Normativa y certificaciones de los EPIs}
    \begin{itemize}
        \item \textbf{Marcado CE} y normas EN correspondientes:
        \begin{itemize}
            \item EN166 para gafas.
            \item EN388 para guantes.
        \end{itemize}
        \item \textbf{Obligaciones del empleador:}
        \begin{itemize}
            \item Proporcionar los EPIs necesarios.
            \item Asegurar su mantenimiento adecuado.
        \end{itemize}
    \end{itemize}
\end{frame}

% Contenidos: Uso y mantenimiento
\begin{frame}{Uso y mantenimiento de los EPIs}
    \begin{itemize}
        \item Colocaci\'on y ajustes adecuados (tallas, cierres, etc.).
        \item Limpieza y revisi\'on peri\'odica para asegurar su eficacia.
        \item Sustituci\'on en caso de deterioro o caducidad.
    \end{itemize}
\end{frame}

% Contenidos: Medidas complementarias
\begin{frame}{Medidas de seguridad complementarias}
    \begin{itemize}
        \item \textbf{Se\~nalizaci\'on} en el puesto de trabajo:
        \begin{itemize}
            \item Zonas de uso obligatorio de gafas, guantes, etc.
        \end{itemize}
        \item \textbf{Protocolos de bloqueo y etiquetado (LOTO)} en equipos el\'ectricos para evitar conexiones accidentales.
        \item \textbf{Checklists diarios o semanales} para revisi\'on de seguridad.
    \end{itemize}
\end{frame}

% Contenidos: Ejemplos en redes
\begin{frame}{Ejemplos concretos en redes y microinform\'atica}
    \begin{itemize}
        \item Uso de pulseras antiest\'aticas para manipular componentes sensibles.
        \item Calzado ESD cuando se trabaja en salas de servidores con moquetas o recubrimientos conductivos.
        \item Gafas de protecci\'on para trabajos con herramientas el\'ectricas de corte o soldadura.
    \end{itemize}
\end{frame}

\section{M5: Gestión Ambiental y Orden en el Lugar de Trabajo}

% Slide 2: Fuentes de Contaminación en Entornos Informáticos

\begin{frame}{Fuentes de Contaminación}
    \begin{itemize}
        \item \textbf{Residuos electrónicos (RAEE)}: placas base, tarjetas, cables.
        \item \textbf{Baterías y pilas usadas}: contienen metales pesados.
        \item \textbf{Tintas y tóner de impresoras}.
        \item \textbf{Productos químicos}: usados para limpieza y desinfección.
    \end{itemize}
    \includegraphics[width=0.6\textwidth]{example_image.jpg}
    \note{
        Ejemplo: Una batería mal desechada puede filtrar cadmio al suelo, contaminando aguas subterráneas. En 2020, una empresa fue multada por un incidente similar en España.
    }
\end{frame}

% Slide 3: Clasificación y Tratamiento de Residuos

\begin{frame}{Clasificación y Tratamiento de Residuos}
    \begin{itemize}
        \item \textbf{Normativa}: residuos peligrosos y no peligrosos.
        \item \textbf{Separación y almacenamiento}: contenedores específicos para RAEE.
        \item \textbf{Procedimientos de recogida selectiva}: empresas autorizadas.
    \end{itemize}
    \includegraphics[width=0.6\textwidth]{recycling_flow.jpg}
    \note{
        Dato interesante: Ecoembes reportó en 2022 un aumento del 15\% en el reciclaje de RAEE en Europa. Este crecimiento muestra cómo la concienciación está teniendo impacto.
    }
\end{frame}

% Slide 4: Fuentes de Contaminación Ambiental

\begin{frame}{Fuentes de Contaminación Ambiental}
    \begin{itemize}
        \item \textbf{Emisión de gases}: productos inflamables o aerosoles.
        \item \textbf{Contaminación de suelos y aguas}: vertidos inadecuados.
        \item \textbf{Impacto energético}: huella de carbono de servidores.
    \end{itemize}
    \note{
        Reflexión: Un servidor promedio en un data center produce aproximadamente 2 toneladas de CO2 al año. Esto equivale a conducir 16.000 km en coche.
    }
\end{frame}

% Slide 5: Orden y Limpieza del Entorno de Trabajo

\begin{frame}{Orden y Limpieza}
    \begin{itemize}
        \item \textbf{Beneficios}:
        \begin{itemize}
            \item Reducción de caídas por cables sueltos.
            \item Localización rápida de herramientas.
        \end{itemize}
        \item \textbf{Principios de las 5S}:
        \begin{enumerate}
            \item Clasificar (Seiri).
            \item Ordenar (Seiton).
            \item Limpiar (Seiso).
            \item Estandarizar (Seiketsu).
            \item Mantener (Shitsuke).
        \end{enumerate}
    \end{itemize}
    \note{
        Ejemplo práctico: Implementar las 5S en un taller de reparación puede reducir un 30\% el tiempo de búsqueda de herramientas.
    }
\end{frame}

\begin{frame}
\frametitle{El origen de las 5S}
\begin{itemize}
    \item La metodología **5S** es parte del \textbf{Sistema de Producción de Toyota (TPS)}.
    \item Nació en Japón como una herramienta clave para optimizar la eficiencia y calidad en los procesos de manufactura.
    \item Su objetivo principal es mantener un entorno de trabajo organizado, seguro y eficiente.
    \item Las **5S** son fundamentales en el modelo de producción \textbf{lean}, ampliamente adoptado en empresas de todo el mundo.
\end{itemize}

\vspace{0.5cm}
\textbf{Frase clave de Toyota:} “Eliminar el desperdicio y fomentar la mejora continua.”
\end{frame}

\section{Clasificar (Seiri)}
\begin{frame}
\frametitle{1. Clasificar (Seiri)}
\begin{itemize}
    \item \textbf{Objetivo:} Identificar y separar los elementos necesarios de los innecesarios.
    \item \textbf{Qué implica:}
        \begin{itemize}
            \item Revisar el material y eliminar lo innecesario.
            \item Priorizar lo esencial para las operaciones.
        \end{itemize}
    \item \textbf{Beneficio:} Reduce el desorden y facilita el acceso a lo necesario.
    \item \textbf{Ejemplo:} Clasificar herramientas en una planta de Toyota, eliminando duplicados o elementos dañados.
\end{itemize}
\end{frame}

\section{Ordenar (Seiton)}
\begin{frame}
\frametitle{2. Ordenar (Seiton)}
\begin{itemize}
    \item \textbf{Objetivo:} Organizar los elementos necesarios en el lugar adecuado.
    \item \textbf{Qué implica:}
        \begin{itemize}
            \item Designar un lugar específico para cada cosa.
            \item Utilizar herramientas como etiquetas o estanterías.
        \end{itemize}
    \item \textbf{Beneficio:} Minimiza el tiempo perdido buscando elementos y mejora la eficiencia.
    \item \textbf{Ejemplo:} Diseñar paneles visuales para herramientas en las líneas de producción de Toyota.
\end{itemize}
\end{frame}

\section{Limpiar (Seiso)}
\begin{frame}
\frametitle{3. Limpiar (Seiso)}
\begin{itemize}
    \item \textbf{Objetivo:} Mantener el lugar de trabajo limpio y en condiciones óptimas.
    \item \textbf{Qué implica:}
        \begin{itemize}
            \item Realizar limpieza regular de equipos y herramientas.
            \item Identificar y corregir fuentes de suciedad.
        \end{itemize}
    \item \textbf{Beneficio:} Mejora la seguridad y el ambiente de trabajo.
    \item \textbf{Ejemplo:} Inspeccionar y limpiar vehículos antes de entregarlos en Toyota.
\end{itemize}
\end{frame}

\section{Estandarizar (Seiketsu) y Mantener (Shitsuke)}
\begin{frame}
\frametitle{4. Estandarizar (Seiketsu) y Mantener (Shitsuke)}
\begin{itemize}
    \item \textbf{Estandarizar (Seiketsu):}
        \begin{itemize}
            \item Crear normas y protocolos claros para mantener las tres primeras S.
            \item Usar listas de verificación y señalización visual.
        \end{itemize}
    \item \textbf{Mantener (Shitsuke):}
        \begin{itemize}
            \item Fomentar hábitos y promover la autorregulación.
            \item Realizar auditorías periódicas y capacitaciones.
        \end{itemize}
    \item \textbf{Beneficio:} Asegura que las mejoras sean sostenibles y parte de la cultura organizacional.
    \item \textbf{Ejemplo:} En Toyota, auditorías regulares aseguran la disciplina en el uso de las 5S.
\end{itemize}
\end{frame}

% Slide 6: Responsabilidad Ambiental y Social

\begin{frame}{Responsabilidad Ambiental y Social}
    \begin{itemize}
        \item \textbf{Buenas prácticas ambientales}: uso de materiales reciclables.
        \item \textbf{Concienciación del personal}: talleres y formación.
        \item \textbf{Certificaciones ambientales}: ISO 14001.
    \end{itemize}
    \note{
        Ejemplo: Una empresa TIC en Barcelona obtuvo la certificación ISO 14001 y redujo su consumo energético en un 20\% en dos años.
    }
\end{frame}





\end{document}
